\documentclass[10pt,a4paper]{article}
\usepackage[margin=2cm]{geometry}
\usepackage[T1]{fontenc}
\usepackage[utf8]{inputenc}
\usepackage{booktabs}
\usepackage{array}
\usepackage{hyperref}
\usepackage{parskip}
\usepackage{enumitem}
\usepackage{xcolor}
\usepackage{titlesec}
\usepackage{multirow}

% Warna header tabel
\definecolor{headerblue}{RGB}{0,70,127}
\definecolor{rowgray}{RGB}{240,240,240}

\titleformat{\section}{\large\bfseries\color{headerblue}}{}{0em}{}[\titlerule]
\titleformat{\subsection}{\normalsize\bfseries}{}{0em}{}

\hypersetup{colorlinks=true, linkcolor=headerblue, urlcolor=headerblue}

\begin{document}

%% ---- HEADER ----
\begin{center}
  {\Large\bfseries Progress Report CP3 --- Desain Eksperimen}\\[4pt]
  {\normalsize Kelompok 08 \quad|\quad 12S4054 Sistem Rekomendasi \quad|\quad Institut Teknologi Del}\\
  {\small Semester Gasal 2026/2027 \quad|\quad Deadline: 7 November 2026}
\end{center}

\vspace{4pt}\hrule\vspace{8pt}

%% ---- 1. RINGKASAN ----
\section{Ringkasan}

CP3 menetapkan desain eksperimen lengkap sebelum eksperimen validasi skala penuh dimulai pada W10--W12.
Protokol evaluasi warm-start, candidate set, leakage rules, hyperparameter search space, dan experiment
matrix telah ditetapkan dan dikodekan dalam repositori GitHub Kelompok 08.

\textbf{Track:} A --- Sparsitas Riwayat User.\\
\textbf{Research Question:} Bagaimana jumlah interaksi yang dimiliki user memengaruhi performa
collaborative, content-based, dan hybrid recommendation?

%% ---- 2. PROTOKOL EVALUASI ----
\section{Protokol Evaluasi (Warm-Start)}

\begin{table}[h!]
\centering
\begin{tabular}{@{}ll@{}}
\toprule
\textbf{Komponen} & \textbf{Nilai} \\
\midrule
Data split & Official temporal split dari paket proyek; tidak di-reshuffle \\
Ukuran training & 427.957 interaksi \\
Ukuran validation & 41.284 (warm: 30.360) \\
Ukuran test & 42.595 (tidak digunakan sebelum W13) \\
Evaluable users & 11.659 pengguna dari \texttt{evaluable\_users\_validation\_warm.csv} \\
Candidate universe & 22.753 item dari \texttt{candidate\_items\_warm.csv} (sama semua model) \\
Ranking per user & Top-10 \\
Metrik wajib & \textbf{Recall@10} dan \textbf{NDCG@10} \\
Item relevan & Rating validation $\geq$ 4 \\
Seen-item masking & Item training dikeluarkan dari kandidat tiap user \\
\bottomrule
\end{tabular}
\end{table}

Model selection dan hyperparameter tuning \textbf{hanya menggunakan validation}. Test set tidak diakses.

%% ---- 3. CANDIDATE SET & LEAKAGE RULES ----
\section{Candidate Set \& Leakage Rules}

\subsection{Candidate Set}
Semua model menggunakan \texttt{candidate\_items\_warm.csv} sebagai candidate universe.
Item yang pernah diinteraksikan user pada training \textbf{dikeluarkan} sebelum rekomendasi dibuat.
Candidate protocol identik untuk seluruh model yang dibandingkan.

\subsection{Leakage Rules}
\begin{enumerate}[leftmargin=*, nosep]
  \item Model hanya dilatih dari data training; validation dan test tidak masuk proses training.
  \item Official temporal split tidak diacak ulang; interaksi validation/test tidak bocor ke training.
  \item User profile dan parameter model dibangun dari training saja. Field agregat
        (\texttt{average\_rating}, \texttt{rating\_number}) tidak digunakan sebagai fitur.
  \item Candidate protocol sama untuk semua model; tidak ada candidate set yang
        menguntungkan model tertentu.
  \item Hyperparameter tuning dan model selection hanya menggunakan validation.
        Test set tidak digunakan sebelum W13.
\end{enumerate}

%% ---- 4. EVALUABLE USERS & ACTIVITY GROUPS ----
\section{Evaluable Users \& Activity Groups}

Pengguna yang masuk evaluasi harus: (a) memiliki riwayat interaksi pada training (warm-start),
dan (b) memiliki ground-truth positif (rating $\geq$ 4) pada validation.

Kelompok aktivitas (Track A) didefinisikan dari jumlah interaksi training menggunakan \texttt{pd.cut}
dengan batas tetap, konsisten untuk semua model (\texttt{make\_user\_groups()} di \texttt{baseline\_common.py}):

\begin{table}[h!]
\centering
\begin{tabular}{@{}lc@{}}
\toprule
\textbf{Kelompok} & \textbf{Jumlah interaksi training} \\
\midrule
Low    & 1 -- 5   \\
Medium & 6 -- 8   \\
High   & $\geq$ 9 \\
\bottomrule
\end{tabular}
\end{table}

%% ---- 5. HYPERPARAMETER SEARCH SPACE ----
\section{Hyperparameter Search Space}

Tuning dilakukan via validation menggunakan \texttt{src/tune\_hyperparameters.py}.
Semua nilai search space didefinisikan di \texttt{src/config.py}.

\begin{table}[h!]
\centering
\begin{tabular}{@{}llll@{}}
\toprule
\textbf{Model} & \textbf{Parameter} & \textbf{Search Space} & \textbf{Metode Seleksi} \\
\midrule
Popularity        & ---                & Tidak ada tuning    & Tetap \\
Matrix Factorization & \texttt{n\_components} & \{32, 64, 128\}     & Recall@10 pada validation \\
TF-IDF CBF        & \texttt{max\_features}  & \{50.000, 100.000\} & Recall@10 pada validation \\
Hybrid            & \texttt{alpha} (bobot MF) & \{0.3, 0.5, 0.7\}  & Recall@10 pada validation \\
\bottomrule
\end{tabular}
\end{table}

Parameter tetap di semua run: \texttt{random\_state=42}, \texttt{n\_iter=7} (MF),
\texttt{ngram\_range=(1,2)} dan \texttt{min\_df=2} (TF-IDF), normalisasi min-max per user (Hybrid).

%% ---- 6. EXPERIMENT MATRIX ----
\section{Experiment Matrix}

\begin{table}[h!]
\centering
\begin{tabular}{@{}llll@{}}
\toprule
\textbf{ID} & \textbf{Model/Variant} & \textbf{Tujuan} & \textbf{Metrik} \\
\midrule
E1      & Popularity                    & Baseline bersama                              & Recall@10, NDCG@10 \\
E2a--c  & MF ($n$ = 32/64/128)          & Tuning collaborative baseline                 & Recall@10, NDCG@10 \\
E3a--b  & TF-IDF CBF (feat = 50k/100k)  & Tuning content-based baseline                 & Recall@10, NDCG@10 \\
E4a--c  & Hybrid ($\alpha$ = 0.3/0.5/0.7) & Tuning hybrid; config dari E2/E3 terbaik    & Recall@10, NDCG@10 \\
E5      & Track-A Analysis              & Analisis per activity group (Low/Medium/High) & Recall@10, NDCG@10 per group \\
\bottomrule
\end{tabular}
\end{table}

File lengkap: \texttt{results/experiment\_matrix.csv}.

%% ---- 7. REPRODUCIBILITY ----
\section{Reproducibility Settings}

\begin{itemize}[leftmargin=*, nosep]
  \item \textbf{Random seed:} \texttt{42} --- dikodekan di \texttt{src/config.py} (\texttt{RANDOM\_SEED}).
  \item \textbf{Config terpusat:} \texttt{src/config.py} --- semua hyperparameter dan search space.
  \item \textbf{Run ID format:} \texttt{<model>\_<param\_key><val>\_seed<seed>}
        (contoh: \texttt{mf\_components64\_seed42}).
  \item \textbf{Experiment log:} setiap run dicatat ke \texttt{results/experiment\_log.csv}
        oleh \texttt{src/experiment\_log.py} beserta timestamp, git commit hash, parameter, dan metrik.
\end{itemize}

%% ---- 8. HASIL AWAL (CP2) ----
\section{Hasil Awal (CP2 --- Belum Dituning)}

\begin{table}[h!]
\centering
\begin{tabular}{@{}lcc@{}}
\toprule
\textbf{Model} & \textbf{Recall@10} & \textbf{NDCG@10} \\
\midrule
Popularity           & 0.0135 & 0.0092 \\
Matrix Factorization & 0.0131 & 0.0079 \\
TF-IDF CBF           & 0.0126 & 0.0073 \\
Hybrid               & \textbf{0.0146} & \textbf{0.0092} \\
\bottomrule
\end{tabular}
\caption{Performa baseline awal pada validation warm (n\_components=32, alpha=0.5, max\_features=100k).}
\end{table}

Hybrid mengungguli semua model. Tuning pada W10--W12 difokuskan pada peningkatan terutama
untuk pengguna kelompok Low yang memiliki riwayat interaksi terbatas.

%% ---- 9. EVIDENCE ----
\section{Evidence GitHub}

File baru/diperbarui pada CP3:
\begin{itemize}[leftmargin=*, nosep]
  \item \texttt{src/config.py} --- konfigurasi dan search space terpusat
  \item \texttt{src/experiment\_log.py} --- logging per run (run ID, timestamp, git commit, metrik)
  \item \texttt{src/tune\_hyperparameters.py} --- skrip tuning otomatis via validation
  \item \texttt{results/experiment\_matrix.csv} --- experiment matrix lengkap (E1--E5)
  \item \texttt{src/mf\_baseline.py}, \texttt{src/tfidf\_cbf\_baseline.py}, \texttt{src/hybrid\_baseline.py}
        --- diperbarui untuk membaca parameter dari \texttt{config.py}
\end{itemize}

\end{document}
